% ==================================================================
% id:       cond-001
% topic:    הסתברות מותנית ואי-תלות
% subtopic: נוסחת ההסתברות השלמה ונוסחת בייס
% solution: full
% verified: no
% ==================================================================
פרופסור אקס רוצה לבחור טוש לקראת ההרצאה שלו. על השולחן במשרדו יש שתי קופסאות טושים: בקופסה א' יש 4 טושים כחולים ו-6 שחורים, ובקופסה ב' יש 3 טושים כחולים ו-7 שחורים. הפרופסור בוחר קופסה באקראי (שתי הקופסאות שוות הסתברות), ומוציא ממנה טוש באקראי.
\begin{solbox}
\unverified
נסמן $A$ = "בחר בקופסה א'" ו-$B$ = "הוציא טוש כחול".
\end{solbox}
\begin{enumerate}
\item מה ההסתברות שהוציא טוש כחול?
\begin{solbox}
\unverified
לפי נוסחת ההסתברות השלמה:
\[
\PP(B)=\PP(B\mid A)\PP(A)+\PP(B\mid\comp A)\PP(\comp A)=\frac4{10}\cdot\frac12+\frac3{10}\cdot\frac12=\frac7{20}=\ans{0.35}.
\]
($\PP(B\mid A)=\frac4{10}$, כי בקופסה א' יש 4 טושים כחולים מתוך 10, ובאותו אופן $\PP(B\mid\comp A)=\frac3{10}$.)

\textbf{דרך נוספת:} מכיוון שהקופסאות שוות הסתברות ובכל אחת 10 טושים, אפשר לחשוב על 20 טושים, מתוכם 7 כחולים, ולכן $\PP(B)=\frac7{20}$.
\end{solbox}
\item \underline{ידוע} שהטוש שהוציא הוא כחול. מה ההסתברות שבחר בקופסה א'?
\begin{solbox}
\unverified
לפי נוסחת בייס:
\[
\PP(A\mid B)=\PP(B\mid A)\cdot\frac{\PP(A)}{\PP(B)}=\frac4{10}\cdot\frac{1/2}{7/20}=\ans{\frac47}.
\]
\end{solbox}
\item האם המאורעות "הפרופסור בחר בקופסה א'" ו"יצא טוש כחול" תלויים או בלתי תלויים?
\begin{solbox}
\unverified
$\PP(B\mid A)=\frac4{10}\neq\frac7{20}=\PP(B)$, ולכן \(\ans{\text{המאורעות $A,B$ תלויים}}\).
לחלופין: $\PP(A\cap B)=\PP(A\mid B)\PP(B)=\frac47\cdot\frac7{20}=\frac15$, ואילו $\PP(A)\PP(B)=\frac12\cdot\frac7{20}=\frac7{40}$.
\end{solbox}
\item אחרי שבוע הפרופסור צריך לבחור שני טושים. הפעם קופסה א' מונחת על השולחן, וקופסה ב' בקצה השני של החדר. מכיוון שהוא מתעצל לפעמים, ההסתברות שיבחר בקופסה א' היא $p$, עבור $p>\frac12$ נתון, וההסתברות שיקום ויבחר בקופסה ב' היא $1-p$.
\begin{enumerate}
\item חשבו את ההסתברות שהפרופסור הוציא שני טושים כחולים (כפונקציה של $p$).
\item ההסתברות שהפרופסור בחר בקופסה ב', בהינתן שהוציא שני טושים כחולים, היא $\frac17$. חשבו את $p$.
\end{enumerate}
\begin{solbox}
\unverified
\begin{enumerate}
\item נסמן $C$ = "הוציא שני טושים כחולים". כאן $\PP(A)=p$ ו-$\PP(\comp A)=1-p$. אם בחר בקופסה א', יש $\binom{10}2=45$ דרכים לבחור שני טושים, מתוכן $\binom42=6$ טובות, ולכן $\PP(C\mid A)=\frac6{45}=\frac2{15}$. באותו אופן $\PP(C\mid\comp A)=\frac{\binom32}{45}=\frac1{15}$. לכן
\[
\PP(C)=\frac2{15}p+\frac1{15}(1-p)=\ans{\frac{1+p}{15}}.
\]
\item נתון $\PP(\comp A\mid C)=\frac17$. לפי בייס:
\[
\frac17=\PP(C\mid\comp A)\cdot\frac{\PP(\comp A)}{\PP(C)}=\frac1{15}\cdot\frac{1-p}{\frac{1+p}{15}}=\frac{1-p}{1+p}
\ \Longrightarrow\ 1+p=7(1-p)\ \Longrightarrow\ \ans{p=\tfrac34}.
\]
\end{enumerate}
\end{solbox}
\end{enumerate}
